\chapter{Estado del arte}
\label{ch:estado-del-arte}
% Síntesis crítica de la literatura (qué se hizo y qué se encontró); las fórmulas van en Fundamentos.

Este capítulo revisa la literatura que sostiene el trabajo. No pretende ser exhaustivo, sino selectivo y crítico: de cada propuesta interesa la idea central, lo que sus autores encontraron y, sobre todo, qué deja sin resolver. Las definiciones formales de cada métrica se posponen al Capítulo~\ref{ch:fundamentos}; aquí se manejan en su versión conceptual para no romper el hilo argumental. La Sección~\ref{sec:eda-prediccion} sitúa el problema general de predecir el resultado de un entrenamiento desde sus primeras épocas, un problema que la búsqueda de arquitecturas neuronales lleva años abordando. La Sección~\ref{sec:sota-alineacion} revisa la familia de métricas de alineación o coherencia direccional, y la Sección~\ref{sec:sota-variabilidad} la de variabilidad estocástica. La Sección~\ref{sec:sota-baselines} introduce los predictores que sirven de referencia y que cualquier métrica de gradiente debe superar para justificar su coste. La Sección~\ref{sec:sota-hueco} cierra el capítulo situando la aportación de este trabajo frente a lo revisado.

\section{Predicción temprana de la eficiencia del entrenamiento}
\label{sec:eda-prediccion}

La idea de anticipar el resultado de un entrenamiento sin llegar a completarlo no es nueva. El campo que más la ha trabajado es la búsqueda de arquitecturas neuronales (NAS, por sus siglas en inglés), donde el coste de entrenar a fondo cada arquitectura candidata es prohibitivo: evaluar un espacio de búsqueda completo puede requerir miles de horas de GPU. La respuesta natural ha sido estimar el rendimiento final de cada candidata a partir de unas pocas épocas de entrenamiento, de modo que el ranking aproximado baste para descartar la mayoría y reservar el cómputo para las prometedoras.

El representante más directo de esta línea es el estimador de velocidad de entrenamiento (TSE), propuesto por Ru et al.~\cite{ru2021speedy}. La premisa que lo sustenta es que las arquitecturas que acabarán generalizando mejor tienden a entrenar más deprisa desde el principio, de modo que una suma del loss de entrenamiento temprano, sin tocar el conjunto de test ni un conjunto de validación, ya ordena bien a las candidatas. Los autores muestran que su estimador correlaciona fuertemente, en términos de Spearman y Kendall, con el ranking real de test, y que lo hace de forma más estable que los llamados proxies de coste cero, que se evalúan en la inicialización sin entrenar nada. Su principal ventaja es el coste: el loss de entrenamiento ya se calcula en cada paso, así que el estimador es prácticamente gratuito.

Esa misma ventaja convierte a TSE en el rival a batir, no en un colaborador. Si la curva de loss temprana ya predice el resultado a coste nulo, cualquier métrica que exija instrumentar el gradiente tiene que demostrar que aporta algo por encima de ella. Este trabajo adopta esa exigencia como criterio central, y por eso TSE reaparece como predictor de referencia en la Sección~\ref{sec:sota-baselines}.

\section{Familia de alineación o coherencia direccional}
\label{sec:sota-alineacion}

La primera de las dos grandes familias agrupa las métricas que miden en qué medida los gradientes de ejemplos distintos apuntan hacia el mismo sitio. La intuición compartida es que, cuando ejemplos distintos producen gradientes alineados, la red está aprendiendo características comunes a muchos de ellos en lugar de memorizarlos uno a uno, y que esa coherencia se asocia con un aprendizaje más rápido y una mejor generalización.

El marco conceptual de la familia es la hipótesis de gradientes coherentes de Chatterjee~\cite{chatterjee2020coherent}. Su punto de partida es una descomposición de la norma del gradiente del minibatch en la energía individual de cada ejemplo más los términos cruzados entre pares: cuando los términos cruzados son positivos, el gradiente agregado se refuerza en las direcciones que muchos ejemplos comparten y se cancela en las idiosincráticas, lo que explicaría por qué el descenso por gradiente estocástico generaliza sin sobreajustar de inmediato. El trabajo es sugerente, pero se queda en lo conceptual: sus medidas operativas exigen partir el conjunto de entrenamiento en ejemplos limpios y corruptos mediante ruido de etiquetas controlado, un montaje que no se traslada a un entrenamiento ordinario.

La versión escalable de esa idea es la m-coherencia de Chatterjee y Zielinski~\cite{chatterjee2020making}. Reduce la coherencia a un único escalar calculable sobre un minibatch, con una interpretación intuitiva: aproximadamente el número medio de ejemplos que se benefician del paso de descenso obtenido a partir de un ejemplo cualquiera. Los autores observan, entrenando una ResNet-18 sobre ImageNet~\cite{russakovsky2015imagenet}, que con etiquetas reales la coherencia crece deprisa, mientras que con etiquetas aleatorias también acaba emergiendo, solo que mucho más despacio. De hecho, que la coherencia surja incluso donde no hay estructura que generalizar es el hallazgo que da título al trabajo; lo que separa el aprendizaje de la memorización no es la presencia de coherencia, sino la velocidad con que aparece. La limitación es de cobertura experimental: la evidencia procede de un único conjunto de datos a gran escala, con un tamaño de batch de 4096 y sin data augmentation, condiciones lejanas a las de un entrenamiento típico, y la métrica se estudia como ventana a la memorización, no como predictor de eficiencia.

La stiffness de Fort et al.~\cite{fort2019stiffness} mira el mismo fenómeno con el coseno entre gradientes de pares de ejemplos, pero desagregándolo por clases. Su matriz de stiffness por clase revela jerarquías semánticas que la red descubre sola, por ejemplo, que mejorar la clasificación de gatos ayuda más a la de perros que a la de camiones, y los autores documentan que la stiffness intra-clase decae hacia cero justo cuando empieza el sobreajuste, lo que la señala como indicador temprano. El inconveniente práctico es estadístico: el estimador se basa en promedios sobre pares de un \emph{conjunto de medición} pequeño, y su variante más informativa exige un \emph{conjunto de medición} mayor para no quedar dominada por el ruido muestral, lo que limita su uso como predictor barato en entrenamientos ordinarios.

La gradient confusion de Sankararaman et al.~\cite{sankararaman2020impact} adopta el punto de vista del peor caso: mide el coseno más negativo entre cualquier par de gradientes, es decir, cuán antagónicos llegan a ser los ejemplos. Su aportación más citada es relacionar esta cantidad con la arquitectura: la confusión crece con la profundidad pero decrece con la anchura, lo que ofrece una explicación cuantitativa de por qué la sobreparametrización facilita el entrenamiento. Los autores muestran además que la \emph{batch normalization} y las conexiones residuales la reducen de forma drástica. Como predictor temprano tiene dos debilidades: al ser un estadístico de extremo (un mínimo) es ruidoso, y su cálculo sobre muchos pares es de los más caros de la familia.

La gradient disparity de Forouzesh y Thiran~\cite{forouzesh2021disparity} sustituye el coseno por la distancia entre los gradientes de dos minibatches independientes, y la propone como señal de parada temprana. Es de las pocas métricas de la familia con evidencia cuantitativa fuerte: los autores obtienen, sobre 220 configuraciones, una correlación de Pearson de 0,957 entre las dos versiones de la disparidad, la que compara minibatches de entrenamiento entre sí y la que compara entrenamiento con validación, lo que justifica usar la primera, que no gasta datos de validación, como sustituto de la segunda. Su punto débil es la comparabilidad: la métrica depende del tamaño de batch (su varianza decrece con él), de modo que solo es comparable entre entrenamientos que lo fijen.

La propuesta más reciente de la familia es la alineación gradiente-peso (GWA) de Hölzl~\cite{hoelzl2025gradient}, que mide el coseno entre el gradiente de cada ejemplo y los pesos del clasificador final. Es la métrica con la correlación publicada más alta del corpus revisado: Pearson de 0,99 y Spearman de 0,97 entre su valor máximo a lo largo del entrenamiento y el accuracy de test. Admite además una forma cerrada que la deja casi gratis cuando se calcula solo sobre la última capa. Sus resultados, sin embargo, se obtienen con arquitecturas modernas de tipo convolucional y transformer y con un protocolo orientado a la robustez frente a ruido de etiquetas, por lo que queda abierto si esa correlación tan alta se sostiene en arquitecturas y condiciones distintas.

\section{Familia de variabilidad estocástica}
\label{sec:sota-variabilidad}

La segunda familia no mira la dirección relativa de los gradientes, sino su dispersión. La intuición es complementaria de la anterior: un gradiente cuya estimación varía mucho de un minibatch a otro da al optimizador una señal ruidosa, y la magnitud de ese ruido respecto a la señal condiciona cuánto se puede avanzar por paso y cuál es el tamaño de batch a partir del cual deja de merecer la pena agrandarlo.

El trabajo de referencia es el modelo empírico del entrenamiento con batch grande de McCandlish et al.~\cite{mccandlish2018empirical}, que introduce la escala de ruido del gradiente (GNS): el cociente entre la traza de la covarianza del gradiente y el cuadrado de su norma. Su aportación es predecir el tamaño de batch crítico, el punto a partir del cual duplicar el batch deja de reducir el número de pasos necesarios, sin necesidad de barrer ese hiperparámetro. La métrica es, por construcción, una predicción de un hiperparámetro de entrenamiento más que de la eficiencia final, y su versión exacta requiere la Hessiana; la forma simplificada que se usa en la práctica asume que la curvatura es proporcional a la identidad, una aproximación cómoda pero no inocua.

Faghri et al.~\cite{faghri2020study} estudian directamente la varianza del gradiente y proponen una varianza normalizada que la hace comparable entre problemas de distinta escala. Su hallazgo más llamativo es contraintuitivo: en CIFAR e ImageNet la varianza normalizada crece durante el entrenamiento en lugar de decrecer, lo que contradice la imagen ingenua de que el gradiente se va estabilizando a medida que la red aprende. Es un trabajo descriptivo más que predictivo: mide y caracteriza la varianza, pero no la usa para anticipar la eficiencia de un entrenamiento.

Liu et al.~\cite{liu2020understanding} cierran la familia con la razón señal-ruido del gradiente por parámetro (GSNR), que es esencialmente el inverso de la varianza normalizada. Su tesis es que un GSNR alto durante el entrenamiento se asocia con un menor gap de generalización, y la respaldan con un marco teórico que relaciona el GSNR con la mejora de generalización por paso, y con correlaciones de Pearson de 0,907 a 0,968 entre ambas cantidades. La evidencia, sin embargo, procede en buena parte de entrenamientos con descenso de gradiente puro y tasas de aprendizaje pequeñas, condiciones que aíslan el fenómeno pero se alejan del entrenamiento estocástico habitual.

Conviene situar aquí, como telón de fondo, dos trabajos que no aportan métricas pero justifican el interés por medir la varianza en lugar de combatirla. El gradiente estocástico de varianza reducida de Johnson y Zhang~\cite{johnson2013accelerating} demuestra que reducir explícitamente la varianza del estimador acelera la convergencia bajo convexidad fuerte, lo que legitima la varianza como cantidad relevante para la velocidad de entrenamiento. Pero Defazio y Bottou~\cite{defazio2019ineffectiveness} muestran, de forma empírica, que esa reducción fracasa en redes profundas modernas: la estructura de suma finita sobre la que se apoya se rompe con la batch normalization, el data augmentation y el dropout, hasta el punto de que el método llega a inyectar más varianza de la que elimina. La conclusión que el presente trabajo extrae de ese contraste es que, en aprendizaje profundo real, la varianza del gradiente tiene más valor como señal diagnóstica que como objetivo a minimizar.

Por último, los propios optimizadores del estudio guardan una relación directa con esta familia. El repaso de Ruder~\cite{ruder2016overview} ofrece la taxonomía de variantes de descenso por gradiente, y Adam~\cite{kingma2015adam} y su antecedente RMSProp~\cite{tieleman2012rmsprop} incorporan en su propia mecánica una razón señal-ruido por parámetro: dividen el gradiente por una estimación de su magnitud reciente, que es justamente lo que mide la familia de variabilidad. Esta coincidencia motiva una decisión de diseño que se justifica en el Capítulo~\ref{ch:metodologia}: medir todas las métricas sobre el gradiente bruto, y nunca sobre el paso ya preacondicionado por el optimizador, para que sean comparables entre SGD y Adam.

\section{Predictores de referencia sin instrumentar el gradiente}
\label{sec:sota-baselines}

El valor de este trabajo depende de una comparación honesta. Medir el gradiente cuesta, así que el rival no es la métrica más simple ni la que cada artículo proponga, sino el mejor predictor que se puede obtener sin instrumentar el gradiente en absoluto, leyendo solo lo que el entrenamiento ya produce. Ese predictor de referencia tiene dos componentes, ambos de coste prácticamente nulo.

El primero es el estimador de velocidad de entrenamiento (TSE) de la Sección~\ref{sec:eda-prediccion}, que resume la curva de loss de entrenamiento temprana. El segundo, y probablemente el más fuerte, es la propia métrica de validación medida en la misma fracción temprana del entrenamiento, es decir, el accuracy o el loss de validación a una fracción dada. Es una señal casi gratuita y muy informativa: si las primeras épocas ya separan las configuraciones buenas de las malas en validación, una métrica de gradiente que no supere esa separación no añade nada que merezca su coste.

Estos dos predictores definen el suelo que las métricas de gradiente deben superar. La comparación no se hace solo en términos de correlación cruda, sino de valor incremental: cuánto añade una métrica una vez descontado lo que la curva de loss y la validación temprana ya predicen. Este es el contraste que el objetivo~\ref{obj:incremental} del Capítulo~\ref{ch:introduccion} señala como decisivo, y su formalización se detalla en el Capítulo~\ref{ch:metodologia}.

\section{Posicionamiento de este trabajo}
\label{sec:sota-hueco}

La revisión anterior deja ver un patrón. Cada métrica se ha propuesto y validado por separado, con un objetivo propio (explicar la generalización, elegir el tamaño de batch o decidir cuándo parar), sobre un conjunto de datos y una arquitectura concretos, y rara vez frente al predictor de referencia barato. La Tabla~\ref{tab:sota-metricas} resume las métricas revisadas en los mismos ejes, lo que hace visible esa dispersión: conviven señales de signo y coste muy distintos, evaluadas en bancos de prueba que no se solapan.

Conviene precisar cómo se lee esa tabla. La columna de señal recoge el sentido que cada trabajo asocia con un entrenamiento más eficiente o una mejor generalización, con $\uparrow$ cuando lo deseable es un valor alto y $\downarrow$ cuando lo es un valor bajo. El coste se expresa en relación con el de un paso de entrenamiento y distingue tres tramos: bajo cuando la sobrecarga queda por debajo del cinco por ciento, medio cuando exige un barrido de gradientes por ejemplo sobre el conjunto de medición, y alto cuando además hay que recorrer todos los pares de ese barrido. La columna de evidencia recoge la correlación más fuerte que cada trabajo publica, y es la que pide más cautela, porque las cifras no son comparables entre sí: solo la de GWA relaciona la métrica con el accuracy de test, mientras que la de gradient disparity compara dos variantes de la propia métrica y la de GSNR contrasta los dos lados de la ecuación teórica que sus autores derivan. Los trabajos que respaldan su métrica con figuras y no con un coeficiente aparecen como evidencia cualitativa. Esa heterogeneidad, más que la magnitud de ningún coeficiente concreto, es lo que justifica someter a todas las métricas a una evaluación común.

\begin{table}[ht]
\centering
\footnotesize
\setlength{\tabcolsep}{3pt}
\caption{Síntesis comparativa de las métricas tempranas revisadas.}
\label{tab:sota-metricas}
\begin{tabular}{>{\raggedright\arraybackslash}p{1.9cm}
                >{\raggedright\arraybackslash}p{2.7cm}
                >{\centering\arraybackslash}p{0.9cm}
                >{\raggedright\arraybackslash}p{2.3cm}
                >{\raggedright\arraybackslash}p{2.4cm}
                >{\raggedright\arraybackslash}p{2.3cm}
                >{\centering\arraybackslash}p{1.0cm}}
\toprule
\textbf{Métrica} & \textbf{Qué mide} & \textbf{Señal} & \textbf{Propósito original} & \textbf{Banco de pruebas} & \textbf{Evidencia publicada} & \textbf{Coste} \\
\midrule
\multicolumn{7}{l}{\itshape Familia de alineación o coherencia direccional} \\
\addlinespace[2pt]
m-coherencia \cite{chatterjee2020making} & Ejemplos que se benefician del paso obtenido de uno solo & $\uparrow$ & Separar aprendizaje de memorización & ImageNet; ResNet-18 & Cualitativa & Medio \\
\addlinespace[2pt]
Stiffness \cite{fort2019stiffness} & Coseno entre gradientes de pares, desglosado por clase & $\uparrow$ & Geometría del aprendizaje y aviso de sobreajuste & MNIST, Fashion-MNIST, CIFAR-10/100; MLP, CNN, ResNet-20 & Cualitativa & Medio \\
\addlinespace[2pt]
Gradient confusion \cite{sankararaman2020impact} & Coseno más negativo entre pares, el peor caso & $\downarrow$ & Efecto de la sobreparametrización & CIFAR-10/100, MNIST; redes residuales anchas & Cualitativa & Alto \\
\addlinespace[2pt]
Gradient disparity \cite{forouzesh2021disparity} & Distancia entre los gradientes de dos minibatches & $\downarrow$ & Criterio de parada temprana & MNIST, CIFAR-10/100; AlexNet, VGG-13, ResNet-18 & $r = 0{,}957$ entre sus dos variantes, sobre 220 configuraciones & Bajo \\
\addlinespace[2pt]
GWA \cite{hoelzl2025gradient} & Coseno entre el gradiente y los pesos del clasificador & $\uparrow$ & Proxy de generalización & CIFAR-10/-N/-C, ImageNet-1k/-C; ConvNeXt-F, ViT-S/16 & $r = 0{,}99$ y $\rho = 0{,}97$ frente al accuracy de test & Bajo \\
\addlinespace[3pt]
\multicolumn{7}{l}{\itshape Familia de variabilidad estocástica} \\
\addlinespace[2pt]
Varianza normalizada \cite{faghri2020study} & Varianza del gradiente relativa a su magnitud & $\downarrow$ & Caracterizar la varianza, sin fin predictivo & MNIST, CIFAR-10/100, ImageNet; MLP, ResNet & Cualitativa & Bajo \\
\addlinespace[2pt]
Escala de ruido (GNS) \cite{mccandlish2018empirical} & Traza de la covarianza sobre la norma del gradiente al cuadrado & $\downarrow$ & Predecir el batch size crítico & MNIST, SVHN, CIFAR-10, ImageNet, texto y aprendizaje por refuerzo & Cualitativa & Medio \\
\addlinespace[2pt]
GSNR \cite{liu2020understanding} & Razón señal-ruido del gradiente por parámetro & $\uparrow$ & Explicar la generalización & MNIST, CIFAR-10; CNN, ResNet-18 & $r = 0{,}907$ a $0{,}968$ entre los lados de su ecuación & Medio \\
\addlinespace[3pt]
\multicolumn{7}{l}{\itshape Predictor de referencia, sin instrumentar el gradiente} \\
\addlinespace[2pt]
TSE \cite{ru2021speedy} & Suma del loss de entrenamiento de las primeras épocas & $\downarrow$ & Ordenar arquitecturas en NAS & NAS-Bench-201/301; CIFAR-10/100, ImageNet-16-120 & Spearman y Kendall altos frente al ranking real & Nulo \\
\bottomrule
\end{tabular}
\end{table}

La aportación no es, por tanto, una métrica nueva, sino una evaluación comparativa y rigurosa de las existentes. El resultado que se busca no es coronar a la que mejor predice, sino a la que mejor predice por unidad de coste, lo que se expresa como un frente de Pareto entre potencia predictiva y coste de cómputo. Conviene acotar también lo que el trabajo no puede afirmar: la evidencia se limita a la clasificación de imágenes, con ImageNet sustituido por su versión reducida, y es de naturaleza correlacional, de modo que las relaciones que se encuentren no autorizan conclusiones causales ni una extrapolación automática a otros dominios. El Capítulo~\ref{ch:metodologia} traduce este posicionamiento en un diseño experimental con hipótesis falsables.
